\documentclass[11pt,openany]{book}

\usepackage{fontspec}
\usepackage{geometry}
\geometry{paper=a4paper, margin=1in}
\usepackage{xcolor}
\usepackage{hyperref}
\usepackage{xurl}
\usepackage{titlesec}
\usepackage{fancyhdr}
\usepackage{microtype}
\usepackage{setspace}
\usepackage[most]{tcolorbox}
\usepackage{enumitem}
\usepackage{tikz}

\setmainfont{TeX Gyre Pagella}
\setsansfont{TeX Gyre Heros}
\setmonofont{DejaVu Sans Mono}
\setstretch{1.08}

\definecolor{ink}{HTML}{1F2530}
\definecolor{ember}{HTML}{8F5A3C}
\definecolor{paper}{HTML}{F8F4EC}
\definecolor{mist}{HTML}{E9EEF2}
\definecolor{shadow}{HTML}{C9B9A6}

\hypersetup{
  colorlinks=true,
  linkcolor=ink,
  urlcolor=ember,
  pdftitle={Quad Tango Muto},
  pdfauthor={LazyingArt LLC}
}

\titleformat{\chapter}[display]
  {\normalfont\sffamily\bfseries\Huge\color{ink}}
  {\filleft\Large\color{ember}\MakeUppercase{Chapter \thechapter}}
  {1ex}
  {\titlerule[1pt]\vspace{1ex}\filright}

\titleformat{\section}
  {\normalfont\sffamily\bfseries\Large\color{ink}}
  {\thesection}
  {0.75em}
  {}

\pagestyle{fancy}
\fancyhf{}
\fancyhead[LE,RO]{\sffamily\small\thepage}
\fancyhead[RE]{\sffamily\small Quad Tango Muto}
\fancyhead[LO]{\sffamily\small LazyingArt LLC}
\setlength{\headheight}{14pt}

\setlist[itemize]{leftmargin=1.4em}
\setlength{\parindent}{0pt}
\setlength{\parskip}{0.75em}

\newtcolorbox{coverbox}{
  enhanced,
  colback=paper,
  colframe=ember,
  boxrule=0.8pt,
  arc=2pt,
  left=14pt,
  right=14pt,
  top=12pt,
  bottom=12pt
}

\newtcolorbox{chapterbox}{
  enhanced,
  colback=mist,
  colframe=ink,
  boxrule=0.6pt,
  arc=2pt,
  left=12pt,
  right=12pt,
  top=10pt,
  bottom=10pt,
  borderline west={3pt}{0pt}{ember}
}

\newtcolorbox{notebox}{
  enhanced,
  colback=paper,
  colframe=shadow,
  boxrule=0.5pt,
  arc=2pt,
  left=12pt,
  right=12pt,
  top=10pt,
  bottom=10pt
}

\newcommand{\chapterprompt}[1]{
\begin{chapterbox}
\textbf{Chapter prompt}\par
#1
\end{chapterbox}
}

\begin{document}
\frontmatter

\begin{titlepage}
\pagecolor{paper}
\color{ink}
\thispagestyle{empty}
\begin{tikzpicture}[remember picture, overlay]
  \draw[color=ember, line width=1pt]
    ([xshift=1.25cm,yshift=-1.25cm]current page.north west)
    rectangle
    ([xshift=-1.25cm,yshift=1.25cm]current page.south east);
  \draw[color=shadow, line width=0.5pt]
    ([xshift=1.55cm,yshift=-1.55cm]current page.north west)
    rectangle
    ([xshift=-1.55cm,yshift=1.55cm]current page.south east);
\end{tikzpicture}

\vspace*{2.2cm}
{\sffamily\Large\color{ember}\textsc{LazyingArt LLC}\par}
\vspace{1.6cm}
{\sffamily\bfseries\fontsize{34}{38}\selectfont Quad Tango Muto\par}
\vspace{0.45cm}
{\Large A manuscript scaffold for later writing\par}
\vspace{1.4cm}
\begin{coverbox}
\textbf{Starter brief}\par
This volume is initialized as a clean, expandable writing surface. It is built to hold the first statement of the book, the shape of its chapters, its fragments, and the notes that will later become a finished manuscript.
\end{coverbox}
\vfill
{\large Author: LazyingArt LLC\par}
{\large April 2026\par}
\end{titlepage}
\nopagecolor

\tableofcontents

\mainmatter

\chapter{Opening position}
\chapterprompt{Write the first true sentence of the book here. Name the tension, motion, or atmosphere that justifies the title \textit{Quad Tango Muto}.}

Use this opening chapter to establish what kind of book this is, what world it enters, and what kind of attention it asks from the reader.

\section{Purpose}

State the reason this manuscript exists.

\section{Immediate questions}

\begin{itemize}
  \item What must the reader understand immediately?
  \item What mystery, claim, or pressure sets the book in motion?
  \item What is the emotional or intellectual temperature of the first pages?
\end{itemize}

\chapter{Foundations}
\chapterprompt{Define the concepts, assumptions, symbols, or conditions the rest of the manuscript depends on.}

Use this chapter to stabilize the vocabulary and architecture of the book.

\section{Terms}

Define the central words and recurring ideas.

\section{Scope}

Clarify what belongs inside the book and what remains outside it.

\chapter{Main development}
\chapterprompt{Build the primary argument, narrative line, investigation, or sequence of transformations here.}

This is the longest structural space in the starter manuscript. Expand it into subchapters later if needed.

\section{Movement}

Describe how the book advances from one idea or state to another.

\section{Evidence}

List the kinds of proof, scenes, references, examples, or formal structures the chapter should use.

\chapter{Counterforce}
\chapterprompt{Record the pressures that resist the book: objections, rival interpretations, ambiguities, or unresolved fractures.}

Use this chapter to keep the manuscript intellectually honest.

\section{Counterarguments}

What intelligent disagreement should the book face directly?

\section{Unresolved edges}

What tensions should remain visible rather than being smoothed away too early?

\chapter{Resolution}
\chapterprompt{Write the final statement that survives the journey of the manuscript. Keep it clear enough to guide revision later.}

Use the closing chapter to show what has changed and what remains.

\section{Last page test}

If the reader remembers only one paragraph, what should it be?

\appendix

\chapter{Working notes}

\begin{notebox}
\textbf{Use this appendix as a live workshop}\par
Store chapter fragments, quotations, title variants, structural sketches, source trails, and alternate endings here until they find a better home in the manuscript.
\end{notebox}

\begin{itemize}
  \item Quotes to revisit
  \item Sources to collect
  \item Chapter fragments
  \item Structural ideas
  \item Discarded passages worth preserving
\end{itemize}

\end{document}
